\chapter{نتیجه‌گیری و پیشنهادها}
\label{chap:conclusion}

در این پژوهش، مسئله ناوبری خودکار یک اسکوتر چهارچرخ در محیط پویا و دارای موانع ایستا و متحرک بررسی شد. هدف اصلی، طراحی و پیاده‌سازی چارچوبی بود که بتواند با استفاده از یادگیری تقویتی عمیق، داده‌های حسگری و لایه ایمنی، رفتار ناوبری به سمت هدف و اجتناب از موانع را در شبیه‌سازی بیاموزد و سپس روی یک سامانه واقعی اجرا شود. برای این منظور، معماری معلم--دانشجو مبتنی بر \lr{LSTM-SAC} توسعه داده شد؛ به‌گونه‌ای که عامل معلم با مشاهده ایده‌آل و عامل دانشجو با مشاهده مبتنی بر حسگر آموزش داده شدند. همچنین، یک سپر ایمنی برای پایش و اصلاح رفتار عامل در شرایط خطرناک طراحی و در مراحل آموزش، ارزیابی و پیاده‌سازی مورد استفاده قرار گرفت.

در فصل‌های پیشین، ابتدا مبانی نظری و روش‌های مرتبط با یادگیری تقویتی، ناوبری خودکار، ادراک عمق و کنترل ایمن بررسی شدند. سپس، ساختار کلی سامانه، محیط شبیه‌سازی، فضای مشاهده و عمل، تابع پاداش، مراحل آموزش، سپر ایمنی، زنجیره ادراک و پیاده‌سازی سخت‌افزاری تشریح گردید. در ادامه، نتایج آموزش، ارزیابی نهایی سیاست‌ها، مقایسه با کنترل‌کننده \lr{MPC}، ارزیابی اجزای سامانه عملی و اجرای واقعی سیاست روی اسکوتر تحلیل شد. در این فصل، جمع‌بندی نهایی نتایج، دستاوردهای اصلی، محدودیت‌ها و مسیرهای پیشنهادی برای ادامه پژوهش ارائه می‌شود.

\section{جمع‌بندی نتایج پژوهش}
\label{sec:final_results_summary}

نتایج آموزش در محیط شبیه‌سازی نشان دادند که عامل معلم و عامل دانشجو هر دو توانسته‌اند الگوی کلی حرکت به سمت هدف و اجتناب از موانع را بیاموزند. عامل معلم به دلیل استفاده از داده‌های ایده‌آل هندسی، رفتار پایدارتر و کم‌نوسان‌تری داشت، در حالی که عامل دانشجو به دلیل استفاده از داده‌های حسگری و مشاهده ادراکی، با عدم‌قطعیت بیشتری مواجه بود. بااین‌حال، روند آموزش عامل دانشجو نشان داد که استفاده از هدایت معلم، آموزش مرحله‌ای و سپر ایمنی توانسته است سیاست ادراکی را به سطح قابل قبولی از عملکرد برساند.

در ارزیابی نهایی شبیه‌سازی، سیاست‌ها در سخت‌ترین سناریوی آزمون بررسی شدند. این سناریو شامل سه مانع ایستا، چهار مانع متحرک و هدف متحرک بود. نتایج نشان دادند که عامل دانشجو با داده‌های حسگری و سپر فعال به نرخ موفقیت 96 درصد و نرخ برخورد صفر دست یافته است. این نتیجه نشان می‌دهد که ترکیب سیاست دانشجو و سپر ایمنی می‌تواند در شرایط پیچیده، رفتار ناوبری محافظت‌شده و قابل قبولی تولید کند. بااین‌حال، زمانی که مداخله فعال سپر حذف شد، نرخ موفقیت دانشجو به 81 درصد کاهش یافت و نرخ برخورد به 17 درصد رسید. بنابراین، سیاست دانشجو به‌تنهایی هنوز ایمنی ذاتی کامل ندارد و عملکرد ایمن آن در سناریوهای دشوار به حضور سپر فعال وابسته است.

عامل معلم با داده‌های ایده‌آل و سپر فعال بهترین عملکرد شبیه‌سازی را نشان داد و به نرخ موفقیت 98 درصد رسید. اما هنگامی که همین عامل با سپر غیرفعال ارزیابی شد، نرخ موفقیت به 94 درصد کاهش یافت و برخورد به 5 درصد رسید. این نتیجه نشان می‌دهد که حتی برای سیاست قوی‌تر معلم نیز وجود سپر ایمنی باعث بهبود عملکرد می شود.

مقایسه با کنترل‌کننده \lr{MPC} نیز یکی از نتایج مهم پژوهش بود. در همان سناریوی نهایی و با داده‌های حسگری، \lr{MPC} بدون سپر تنها به نرخ موفقیت 30 درصد رسید و نرخ برخورد آن 70 درصد بود. با فعال‌سازی سپر، نرخ موفقیت \lr{MPC} به 50 درصد افزایش یافت، اما نرخ برخورد همچنان 48 درصد باقی ماند. این نتایج نشان دادند که در محیطی پویا، دارای موانع متحرک، هدف متحرک و داده‌های حسگری نامطمئن، کنترل‌کننده پیش‌بین کلاسیک با افق محدود و مدل ساده‌شده نتوانسته است به سطح عملکرد عامل‌های یادگیرنده برسد. از سوی دیگر، هزینه محاسباتی \lr{MPC} نیز نسبت به سیاست عصبی بیشتر بود. بنابراین، سیاست یادگیری‌شده پس از آموزش، هم از نظر کیفیت تصمیم‌گیری در سناریوی موردنظر و هم از نظر هزینه تولید فرمان، مزیت عملی نشان داد.

در بخش پیاده‌سازی عملی، ابتدا اجزای اصلی سامانه واقعی ارزیابی شدند. آزمون فاصله‌سنجی استریو در حالت ایستا نشان داد که زنجیره عمق در شرایط کنترل‌شده دارای دقت قابل قبول است. میانگین قدرمطلق خطا برابر با $0.066~m$، ریشه میانگین مربعات خطا برابر با $0.0812~m$ و میانگین قدرمطلق خطای نسبی برابر با $3.50\%$ به دست آمد. این نتایج نشان دادند که فاصله‌سنجی استریو در حالت ایستا برای تشکیل مشاهده عامل قابل استفاده است.

بااین‌حال، ارزیابی زنجیره تخمین فاصله نسبی در اجرای واقعی اسکوتر نشان داد که خطا در شرایط حرکتی افزایش می‌یابد. این افزایش خطا تنها به دوربین استریو مربوط نیست، بلکه حاصل ترکیب خطای عمق، خطای تشخیص ناحیه هدف یا مانع، لرزش دوربین، خطای تخمین موقعیت اسکوتر از حسگر اینرسی و انکودرها، تأخیر پردازش و خطای تبدیل مختصات است. بنابراین، در تحلیل اجرای واقعی نباید دقت ادراک را با دقت حالت ایستای آزمایشگاهی یکسان در نظر گرفت.

ارزیابی حسگر اینرسی نشان داد که خروجی \lr{Euler yaw} پس از صفرسازی اولیه و بازگشایی زاویه، برای تخمین جهت‌گیری اسکوتر مناسب‌تر از خروجی‌های دیگر است. همچنین، ارزیابی پاسخ حلقه‌بسته موتورهای محرک نشان داد که سامانه حرکتی توانایی دنبال‌کردن روند کلی فرمان‌های سرعت را دارد، اما پاسخ آن ایده‌آل نیست. وجود ناحیه مرده، اصطکاک سکون، بیش‌جهش، تأخیر و ناهمسانی میان چهار موتور نشان داد که بخشی از اختلاف میان شبیه‌سازی و اجرای واقعی به محدودیت‌های عملگرها و کنترل سرعت مربوط است.

در نهایت، سیاست نهایی روی اسکوتر واقعی در قالب 20 اپیزود مستقل اجرا شد. از این میان، 16 اپیزود با موفقیت به هدف رسیدند، 3 اپیزود به اتمام زمان منتهی شدند و 1 اپیزود با برخورد پایان یافت. این نتیجه، معادل نرخ موفقیت 80 درصد، نرخ اتمام زمان 15 درصد و نرخ برخورد 5 درصد است. با توجه به پیچیدگی سامانه واقعی، محدودیت تعداد آزمایش‌ها، نویز حسگرها، خطای تخمین وضعیت و غیرایده‌آل بودن موتورها، این نتیجه نشان می‌دهد که زنجیره کامل ادراک، تصمیم‌گیری و فرمان‌دهی در عمل قابل اجرا بوده است. بااین‌حال، وجود اپیزودهای ناموفق نشان می‌دهد که انتقال از شبیه‌سازی به واقعیت هنوز کامل نیست و نیاز به بهبودهای بیشتر دارد.

\section{دستاوردهای اصلی پژوهش}
\label{sec:main_contributions}

نخستین دستاورد این پژوهش، طراحی یک چارچوب معلم--دانشجو برای ناوبری اسکوتر خودران در محیط پویا است. در این چارچوب، عامل معلم با مشاهده ایده‌آل آموزش دید و عامل دانشجو با مشاهده مبتنی بر حسگرها به سمت رفتار قابل انتقال به سامانه واقعی هدایت شد. این ساختار باعث شد از مزیت مشاهده دقیق در آموزش اولیه استفاده شود، در حالی که سیاست نهایی به داده‌های قابل دسترس در پیاده‌سازی عملی نزدیک‌تر باشد.

دستاورد دوم، استفاده از آموزش مرحله‌ای برای افزایش تدریجی دشواری محیط بود. محیط آموزش از سناریوهای ساده‌تر آغاز شد و به‌تدریج به سناریوی نهایی شامل چند مانع ایستا، چند مانع متحرک و هدف متحرک سریع رسید. این روش باعث شد عامل‌ها بتوانند ابتدا رفتار پایه رسیدن به هدف را بیاموزند و سپس با شرایط پیچیده‌تر شامل موانع متحرک و داده‌های ادراکی مواجه شوند.

دستاورد سوم، طراحی و به‌کارگیری سپر ایمنی در کنار سیاست یادگیری‌شده است. نتایج نشان دادند که سپر ایمنی نقش مهمی در کاهش برخورد، به‌ویژه برای عامل دانشجو، دارد. مقایسه آزمون‌های با سپر فعال و بدون مداخله فعال سپر نشان داد که این لایه ایمنی یکی از اجزای کلیدی برای تبدیل سیاست یادگیری‌شده به یک سامانه قابل استفاده در سناریوهای دشوار است.

دستاورد چهارم، مقایسه روش پیشنهادی با خط مبنای \lr{MPC} در سناریوی یکسان بود. نتایج نشان دادند که سیاست یادگیری‌شده در محیط پویا و مبتنی بر داده حسگری عملکرد بهتری از \lr{MPC} دارد. این مقایسه ارزش روش پیشنهادی را روشن‌تر کرد، زیرا نشان داد عامل یادگیرنده تنها در مقایسه با نسخه‌های دیگر خودش موفق نیست، بلکه نسبت به یک روش کلاسیک کنترل پیش‌بین نیز در سناریوی موردنظر برتری دارد.

دستاورد پنجم، پیاده‌سازی زنجیره کامل سامانه روی اسکوتر واقعی است. در این پژوهش، اجزای مختلف شامل دوربین استریو، حسگر اینرسی، انکودرها، کنترل موتور، تخمین فاصله نسبی، اجرای سیاست و ارسال فرمان به سامانه حرکتی به‌صورت یکپارچه پیاده‌سازی و ارزیابی شدند. اجرای 20 اپیزود واقعی نشان داد که چارچوب پیشنهادی فقط در سطح شبیه‌سازی باقی نمانده و قابلیت اجرای عملی اولیه را نیز دارد.

\section{پاسخ به اهداف پژوهش}
\label{sec:objectives_response}

هدف نخست پژوهش، طراحی سیاستی برای ناوبری اسکوتر در محیط پویا و دارای موانع بود. نتایج آموزش و ارزیابی نهایی نشان دادند که این هدف در محیط شبیه‌سازی تا حد قابل قبولی محقق شده است. عامل دانشجو با داده‌های حسگری و سپر فعال توانست در دشوارترین سناریوی آزمون به نرخ موفقیت بالا و برخورد صفر برسد.

هدف دوم، کاهش فاصله میان آموزش شبیه‌سازی و پیاده‌سازی عملی بود. برای این منظور، عامل دانشجو با داده‌های حسگری آموزش داده شد و زنجیره ادراک واقعی شامل فاصله‌سنجی استریو و تخمین وضعیت حرکتی توسعه یافت. نتایج اجرای واقعی نشان دادند که سیاست آموزش‌دیده توانسته است در اغلب اپیزودهای عملی مأموریت را با موفقیت انجام دهد. بنابراین، این هدف به‌صورت اولیه محقق شده است، هرچند شکاف شبیه‌سازی تا واقعیت همچنان باقی است.

هدف سوم، افزایش ایمنی سامانه در حضور عدم‌قطعیت بود. نتایج نشان دادند که سپر ایمنی در کاهش برخورد نقش مؤثری دارد. برای عامل دانشجو، حذف مداخله فعال سپر باعث افزایش نرخ برخورد شد و در مقابل، فعال‌بودن آن برخورد را در آزمون نهایی شبیه‌سازی به صفر رساند. بنابراین، سپر ایمنی توانسته است به‌عنوان یک لایه حفاظتی مؤثر عمل کند.

هدف چهارم، ارزیابی عملی سامانه روی پلتفرم واقعی بود. این هدف با اجرای 20 اپیزود روی اسکوتر واقعی دنبال شد. نتایج به‌دست‌آمده نشان دادند که سامانه در 80 درصد آزمایش‌ها موفق بوده و تنها در 5 درصد آزمایش‌ها برخورد رخ داده است. بنابراین، امکان اجرای عملی چارچوب پیشنهادی تأیید شد، هرچند برای نتیجه‌گیری آماری گسترده‌تر، تعداد آزمایش‌های بیشتری لازم است.

\section{محدودیت‌های پژوهش}
\label{sec:research_limitations}

نخستین محدودیت پژوهش، وابستگی عملکرد نهایی به سپر ایمنی است. اگرچه عامل دانشجو در حالت سپر فعال عملکرد بسیار خوبی داشته است، اما حذف مداخله فعال سپر باعث افزایش برخورد شد. بنابراین، سیاست دانشجو هنوز به‌تنهایی به سطح ایمنی کامل نرسیده و استفاده ایمن از آن در سناریوهای دشوار نیازمند حضور لایه ایمنی فعال است.

محدودیت دوم، وجود شکاف میان شبیه‌سازی و واقعیت است. در شبیه‌سازی، حتی در حالت استفاده از داده‌های حسگری، بسیاری از شرایط محیط واقعی مانند لرزش دوربین، تغییرات نور، تأخیر پردازش، خطای کالیبراسیون، لغزش چرخ‌ها، ناصافی سطح و ناهمسانی موتورها به‌طور کامل مدل نشده‌اند. این عوامل در اجرای واقعی باعث افت عملکرد نسبت به نتایج شبیه‌سازی می‌شوند.

محدودیت سوم، دقت محدود زنجیره ادراک در شرایط حرکتی است. آزمون‌های فاصله‌سنجی نشان دادند که دقت استریو در حالت ایستا مناسب است، اما در حرکت واقعی خطای تخمین فاصله افزایش می‌یابد. از آنجا که سیاست دانشجو بر پایه همین داده‌های ادراکی تصمیم‌گیری می‌کند، کیفیت ادراک یکی از عوامل تعیین‌کننده در عملکرد نهایی سامانه است.

محدودیت چهارم، غیرایده‌آل بودن سامانه حرکتی است. پاسخ موتورها در سرعت‌های پایین تحت تأثیر ناحیه مرده، اصطکاک سکون و کوانتیزه‌شدن اندازه‌گیری قرار دارد. همچنین، چهار موتور رفتار کاملاً یکسانی ندارند. این ناهمسانی می‌تواند باعث ایجاد خطای مسیر، چرخش ناخواسته و نیاز به اصلاحات مداوم در حرکت شود.

محدودیت پنجم، تعداد محدود آزمایش‌های واقعی است. اگرچه 20 اپیزود واقعی برای اعتبارسنجی اولیه ارزشمند است، اما برای نتیجه‌گیری آماری قوی درباره عملکرد سامانه در همه شرایط محیطی کافی نیست. آزمایش‌های بیشتر در محیط‌های متفاوت، با آرایش‌های گوناگون موانع و شرایط نوری مختلف برای ارزیابی جامع‌تر لازم است.

محدودیت ششم، ساده‌سازی برخی جنبه‌های محیط و هدف است. در این پژوهش، هدف و موانع در چارچوب سناریوهای مشخص تعریف شدند. در محیط‌های واقعی پیچیده‌تر، ممکن است موانع رفتارهای غیرقابل پیش‌بینی‌تری داشته باشند، مسیرها محدودتر باشند یا شرایط تشخیص هدف دشوارتر شود. بنابراین، تعمیم مستقیم نتایج به همه محیط‌های شهری یا نیمه‌ساختاریافته نیازمند آزمایش‌های بیشتر است.

\section{پیشنهادها برای کارهای آینده}
\label{sec:future_work}

یکی از مسیرهای مهم ادامه پژوهش، بهبود زنجیره ادراک است. استفاده از روش‌های پیشرفته‌تر تخمین عمق، ترکیب استریو با حسگرهای مکمل، بهبود تشخیص هدف و مانع، و طراحی فیلترهای مقاوم‌تر برای کاهش نویز و تأخیر می‌تواند کیفیت مشاهده عامل را افزایش دهد. از آنجا که نتایج نشان دادند داده‌های حسگری یکی از گلوگاه‌های اصلی سامانه هستند، بهبود ادراک می‌تواند مستقیماً به افزایش ایمنی و موفقیت منجر شود.

پیشنهاد دوم، استفاده از روش‌های پیشرفته‌تر انتقال از شبیه‌سازی به واقعیت است. روش‌هایی مانند تصادفی‌سازی دامنه، مدل‌سازی نویز حسگرها، شبیه‌سازی تأخیر پردازش، افزودن مدل دقیق‌تر عملگرها و آموزش با اغتشاش‌های دینامیکی می‌توانند مقاومت سیاست را نسبت به شرایط واقعی افزایش دهند. در این پژوهش، بخشی از شکاف شبیه‌سازی و واقعیت مشاهده شد و کاهش این شکاف می‌تواند یکی از محورهای اصلی ادامه کار باشد.

پیشنهاد سوم، بهبود سیاست دانشجو از نظر ایمنی ذاتی است. نتایج نشان دادند که دانشجو در حالت بدون مداخله فعال سپر هنوز برخوردهای قابل توجهی دارد. بنابراین، می‌توان در پژوهش‌های آینده از روش‌هایی مانند یادگیری تقویتی ایمن، قیود ایمنی در تابع هدف، آموزش با نمونه‌های نزدیک به خطر، یادگیری از مداخلات سپر و ترکیب سیاست با تخمین ریسک استفاده کرد تا وابستگی سیاست به سپر کاهش یابد.

پیشنهاد چهارم، توسعه سپر ایمنی تطبیقی است. در نسخه فعلی، سپر بر اساس معیارهای مشخصی مانند فاصله، زمان تا برخورد و سطح خطر عمل می‌کند. در آینده می‌توان سپری طراحی کرد که آستانه‌های آن بر اساس سرعت، تراکم موانع، کیفیت ادراک، عدم‌قطعیت تخمین و شرایط دینامیکی سامانه به‌صورت تطبیقی تغییر کند. چنین سپری می‌تواند در شرایط ایمن کمتر مداخله کند و در شرایط پرخطر محافظه‌کارتر باشد.

پیشنهاد پنجم، بهبود کنترل سطح پایین موتورها است. نتایج آزمون پاسخ موتور نشان دادند که وجود ناحیه مرده، ناهمسانی موتورها و بیش‌جهش در برخی سرعت‌ها می‌تواند بر حرکت واقعی اسکوتر اثر بگذارد. تنظیم مستقل کنترل‌کننده برای هر موتور، استفاده از کنترلرهای مقاوم‌تر، جبران دقیق‌تر اصطکاک، و مدل‌سازی بهتر دینامیک چرخ‌ها می‌تواند اجرای فرمان‌های سیاست را دقیق‌تر کند.

پیشنهاد ششم، افزایش تعداد و تنوع آزمایش‌های واقعی است. برای ارزیابی کامل‌تر، لازم است آزمون‌های بیشتری در محیط‌های متفاوت، با نورپردازی مختلف، سطح‌های حرکتی گوناگون، آرایش‌های متنوع موانع و هدف‌های با رفتار متفاوت انجام شود. همچنین، ثبت و تحلیل دقیق‌تر مسیر حرکت، سرعت، خطای ادراک، مداخلات سپر و علت شکست هر اپیزود می‌تواند به شناخت بهتر نقاط ضعف سامانه کمک کند.

پیشنهاد هفتم، بررسی ترکیب سیاست یادگیری‌شده با کنترل‌کننده‌های مدل‌محور است. اگرچه \lr{MPC} در آزمون‌های این پژوهش عملکرد ضعیف‌تری نسبت به سیاست یادگیری‌شده داشت، اما روش‌های مدل‌محور همچنان می‌توانند در قالب لایه اصلاح‌کننده، پیش‌بینی‌کننده خطر یا تولیدکننده قید برای سیاست یادگیری‌شده مفید باشند. ترکیب یادگیری تقویتی با کنترل پیش‌بین یا کنترل ایمن می‌تواند به سامانه‌ای منجر شود که هم سرعت تصمیم‌گیری سیاست عصبی را داشته باشد و هم از ساختار تحلیلی کنترل مدل‌محور بهره ببرد.

پیشنهاد هشتم، توسعه معیارهای ارزیابی فراتر از موفقیت و برخورد است. در ادامه پژوهش می‌توان شاخص‌هایی مانند راحتی حرکت، مصرف انرژی، میزان نوسان فرمان، فاصله ایمن متوسط، تعداد مداخلات سپر، سطح اعتماد به ادراک و میزان تطابق مسیر واقعی با مسیر مطلوب را نیز به‌صورت دقیق‌تر بررسی کرد. چنین شاخص‌هایی برای کاربردهای واقعی، به‌ویژه در سامانه‌های حمل‌ونقل شخصی، اهمیت زیادی دارند.

\section{نتیجه‌گیری نهایی}
\label{sec:final_conclusion}

این پژوهش نشان داد که استفاده از چارچوب معلم--دانشجو مبتنی بر یادگیری تقویتی عمیق، همراه با سپر ایمنی و زنجیره ادراک حسگری، می‌تواند راهکاری عملی برای ناوبری اسکوتر خودران در محیط‌های پویا فراهم کند. عامل دانشجو توانست با داده‌های حسگری و سپر فعال در سخت‌ترین سناریوی شبیه‌سازی عملکردی نزدیک به عامل معلم به دست آورد و نسبت به خط مبنای \lr{MPC} عملکرد بهتری نشان دهد. همچنین، اجرای واقعی روی اسکوتر نشان داد که زنجیره کامل سامانه در عمل قابل استفاده است و می‌تواند در اغلب اپیزودها مأموریت را با موفقیت انجام دهد.

بااین‌حال، نتایج پژوهش به‌روشنی نشان دادند که عملکرد موفق سامانه حاصل ترکیب چند جزء است: سیاست یادگیری‌شده، مشاهده حسگری، تخمین وضعیت، کنترل موتور و سپر ایمنی. هیچ‌یک از این اجزا به‌تنهایی برای تضمین عملکرد ایمن کافی نیستند. به‌ویژه، ایمنی عامل همچنان به سپر فعال وابسته است و خطاهای ادراکی و حرکتی در اجرای واقعی می‌توانند عملکرد را کاهش دهند.

بنابراین، نتیجه نهایی این پژوهش آن است که چارچوب پیشنهادی یک گام مؤثر به سمت ناوبری خودکار و ایمن اسکوتر در محیط‌های پویا است، اما برای دستیابی به عملکرد پایدارتر و قابل اعتمادتر در کاربردهای واقعی، بهبود بیشتر در ادراک، کنترل حرکتی، آموزش مقاوم نسبت به عدم‌قطعیت و طراحی سپر ایمنی تطبیقی ضروری است.